\documentclass[11pt,a4paper]{article}
\usepackage[T1]{fontenc}
\usepackage[utf8]{inputenc}
\usepackage[margin=27mm]{geometry}
\setlength{\emergencystretch}{2em}
\usepackage{amsmath}
\usepackage{amssymb}
\usepackage{amsthm}
\newtheorem{proposition}{Proposition}[section]
\title{Graded carriers and physical readouts: conditions for mass and geometric interpretation}
\author{Marc Brendecke\\\small Independent Researcher, Bochum, Germany}
\date{Working paper 5.0.0-draft\\\small 11 November 2026; not peer reviewed}
\begin{document}
\maketitle
\begin{abstract}
We examine the conditions under which a common graded state space can
support particle and geometric readouts. The Quantum Sphaera construction
is formulated here as a set of conditional interfaces rather than a
completed identification with physical fields. For a real orthogonal
involution, the associated compensator gives an exact projection, but
does not determine a three-generation mass spectrum. We derive the
additional equal-norm condition for the Koide root ratio and express it
as an orthogonality condition in a declared HC representation. Direct
odd-input and cyclic-equivariant readouts exhibit obstructions, while
positive quadratic readouts require a restricted state domain. A
projection of the dynamics supplies a separate autonomy condition for
observable variables. Together these results isolate the operator,
metric and state-selection data needed to turn an algebraic framework
into a physical model. No lepton masses, absolute mass scale or
gravitational field equations are predicted by the present construction.
\end{abstract}
\section{\raggedright Introduction}
Attempts to describe different physical sectors on a common mathematical
carrier face a recurring difficulty: a shared state space does not by
itself specify the observables that distinguish those sectors. The same
algebra can support different representations, couplings and measurement
maps. A physical interpretation therefore requires more than a matching
number of components or a symmetry with the expected name.

The Quantum Sphaera Companion develops a common-carrier proposal using
biquaternionic and division-algebra structures, an involution and a
compensating projection \cite{QSC411}. This paper extracts a connected
question from that broader programme: under what conditions can the
graded structure be transferred to particle, mass and geometric
readouts without inserting the desired observables into the transfer?
The emphasis is on explicit conditions and consistency constraints.

The mass sector makes this distinction particularly sharp. A relation
among three positive roots can be represented algebraically in many ways,
but a representation selected after the roots are specified is not a
mechanism for their values. We use the Koide ratio as a test of this
interface, derive its missing norm condition, and compare linear and
quadratic readouts. The same distinction reappears in projected dynamics:
an equation on the full carrier need not induce an autonomous equation
for the variables identified as observables.

The scope is consequently conditional. The algebraic projection and
the stated comparison results are available; the identification of a
native carrier, its couplings and its physical observables remains to be
completed. This distinction also governs the particle assignments and
geometric constructions discussed below.
\subsection{Observable equivalence and sufficient state variables}

A common carrier is useful only together with an observable equivalence
relation. Let $\mathcal X$ be an admitted state domain, $\Phi_t$ its
full evolution, and $\pi:\mathcal X\to\mathcal Y$ a readout.
Two states are equivalent for this readout when $\pi(x)=\pi(x')$.
An induced evolution $\phi_t$ on the image exists precisely when
\begin{equation}
 \pi(x)=\pi(x')\quad\Longrightarrow\quad
 \pi(\Phi_t x)=\pi(\Phi_t x')
\end{equation}
for all admitted states and times. Necessity follows by applying
$\phi_t$ to the common initial readout. For sufficiency, define
$\phi_t(y)=\pi(\Phi_t x)$ using any representative with $\pi(x)=y$;
the condition makes the definition independent of the representative.
For a full flow preserving the domain, its composition law descends
as well. Existence, regularity, and locality of a field flow remain
separate analytic questions.

This elementary criterion connects two seemingly different issues.
A compressed field register can omit orientation needed for a Lorentz
transformation. A projected state can omit an amplitude needed for
the derivative of its phase. In both cases equivalence classes of
the original readout are split by the proposed operation. The remedy
must supply the missing information, restrict the admitted domain,
or describe a family of possible outputs. It cannot be supplied by
renaming the original register.

\subsection{Three tests of a proposed physical interface}

We organize the investigation around three questions. First, what
is represented? This requires a carrier, scalar field, adjoint and
operator action. The distinction between an involution and an internal
complex structure is already a necessary test at this level. Second,
which operation selects the physical states? Closure, positive
spectral selection and a generation seed are different restrictions.
Their intersection may have smaller dimension than any one separately.
Third, which measured quantity is read out? An expectation, a transition
norm, a matrix eigenvalue and a pole of a dressed propagator need not
give the same mass law.

The answers can be checked before a detailed phenomenological fit.
Rank rules out three positive eigenmasses from a one-dimensional
transition output. Positivity bounds exclude a globally positive
linear-density realization at the intended root radius. A hidden-lift
counterexample excludes a universal phase-only evolution for a given
projection. These are useful negative results because their hypotheses
identify which part of a model would have to change. They do not
exclude every representation on the common carrier.

The presentation follows these interfaces rather than the historical
order of research rounds. After the carrier and its comparison metrics,
an electromagnetic example explains the assertion and execution
language. Particle representations are separated from generation
seeds and mixing. The mass section then follows a transition operator
through its Gram matrix, a gapped effective comparison, and the positive
root constraints. The geometric section examines projected evolution,
action reduction, coframes and finite compatibility. Finally, a zeta
comparison tests whether arithmetic and uncertainty can supply an
additional selection operation. Original supporting mass arguments are
retained in Appendix~A; the full Companion \cite{QSC411} supplies the
broader research setting and alternative routes.
\section{\raggedright A graded carrier and its interpretation}
Let $H$ be a real Hilbert space and let $J:H\to H$ be an orthogonal
involution. Since $J^{-1}=J=J^*$, the operators
\begin{equation}
 C=\frac{I+J}{2},\qquad P_-=\frac{I-J}{2}
\end{equation}
are complementary orthogonal projections. Indeed,
$C^2=(I+2J+J^2)/4=C$, and similarly $P_-^2=P_-$, while
$CP_-=0$ and $C+P_-=I$. We refer to $C$ as the HC compensator.
It selects the $J$-even part of a state; $P_-$ selects the odd part.

This grading is the minimal structure needed for the subsequent tests.
It does not identify either range with a particle species. In the
Companion, a biquaternionic carrier and extensions involving
$\mathbb C\otimes\mathbb H\otimes\mathbb O$ motivate proposed
particle and frame assignments. Division algebras provide a mathematical
setting for such representation questions \cite{Baez2002}. A positive
Hilbert metric used in a readout calculation must nevertheless be
distinguished from a Lorentzian spacetime metric.

The scalar field also matters. An antiunitary involution on a complex
Hilbert space is generally only real-linear when averaged with the
identity. Here $J$ is explicitly real orthogonal; when complex amplitudes
are used, we mean its declared complex-linear extension. No equivalence
with an inherited antiunitary convention is presumed. Likewise, a change
between the proposed matter, light and tachyon frames requires an
explicit map, not merely a relabelling of components.
\subsection{The algebraic carrier and its two conjugations}
A concrete carrier is needed before the word ``sector'' acquires a
mathematical meaning. Write a biquaternion as
\begin{equation}
 q=q_0+z_1e_1+z_2e_2+z_3e_3,\qquad q_0,z_k\in\mathbb C,
 \qquad e_ie_j=-\delta_{ij}+\epsilon_{ijk}e_k.
\end{equation}
The scalar imaginary unit commutes with the quaternion units. This is
different from choosing a quaternion direction as an internal complex
structure. In particular, a statement formulated with the commuting
scalar unit cannot automatically be transferred to left multiplication
by one of the $e_k$. The representation and the relevant adjoint must
be specified for each transfer.

There are also two distinct conjugations. Define quaternion conjugation
by $\widetilde q=q_0-\sum_kz_ke_k$, leaving the complex coefficients
unchanged, and coefficient conjugation by
$\overline q=\overline{q_0}+\sum_k\overline{z_k}e_k$. Then
\begin{equation}
 q\widetilde q=q_0^2+z_1^2+z_2^2+z_3^2,
\end{equation}
whereas the positive coefficient norm is
\begin{equation}
 \|q\|_E^2=|q_0|^2+|z_1|^2+|z_2|^2+|z_3|^2.
\end{equation}
The first expression is a complex quadratic form; the second is a
positive Hermitian norm. Using the same word ``norm'' for both can
conceal a change of mathematical problem. A positivity argument about
principal mass roots uses a positive comparison. It cannot use a
complex quadratic invariant as though it were that comparison.

One familiar realization sends $e_k$ to $-i\sigma_k$, with the Pauli
matrices $\sigma_k$. The biquaternion becomes a complex two-by-two
matrix, and its determinant is $q_0^2+\sum_kz_k^2$. On the Hermitian
slice $X=tI+\boldsymbol x\cdot\boldsymbol\sigma$ this gives
\begin{equation}
 \det X=t^2-|\boldsymbol x|^2,\qquad
 X\longmapsto AXA^*,\quad A\in SL(2,\mathbb C).
\end{equation}
The determinant is preserved. With the $(-,+,+,+)$ spacetime
convention used in the geometric comparisons, the Lorentz quadratic
form on this slice is $-\det X$. This conventional sign does not
turn the coefficient norm into a Lorentz-invariant Hilbert norm.
For example, $A=\operatorname{diag}(e^{\chi/2},e^{-\chi/2})$
is a boost representative and is not a unitary matrix when
$\chi\ne0$. A similarity, a physical observer transformation and
a unitary change of a positive comparison basis must consequently be
tracked as different operations.

\subsection{What the sign of a field invariant actually distinguishes}
An electromagnetic comparison gives a useful, independently interpretable
instance of the complex vector part. In units with $c=1$, put
$z=\boldsymbol E+i\boldsymbol B$. Then
\begin{equation}
 s=z\cdot z=|\boldsymbol E|^2-|\boldsymbol B|^2
                 +2i\boldsymbol E\cdot\boldsymbol B.
\end{equation}
This identity distinguishes the two real electromagnetic invariants.
Only on a real-$s$ slice, for example
$\boldsymbol E\cdot\boldsymbol B=0$, is a classification by
$s>0$, $s<0$ and $s=0$ directly meaningful. The first two cases are
electric- and magnetic-dominated respectively; the last is a null
invariant when the field is nonzero. Extending a sign label to a
general complex $s$ would require another declared rule.

These algebraic categories are useful for the Companion's proposed
matter, tachyon and photon readouts, but their physical identification
requires dynamics and observables. A null electromagnetic invariant,
for instance, does not by itself specify a quantum photon state.
Nor does a common carrier identify different invariant strata as the
same physical state. The distinction is operational: one must specify
which preparations, transformations and measurements are being compared.

The HC construction is formulated on a positive Hilbert realization
precisely so that its projection identities can be tested independently
of these interpretations. In such a realization every $f$ decomposes
uniquely as
\begin{equation}
 f=f_++f_-,\qquad f_+=Cf,\quad f_-=P_-f,\qquad
 \|f\|^2=\|f_+\|^2+\|f_-\|^2.
\end{equation}
For a normalized state its involution expectation is
\begin{equation}
 \langle f,Jf\rangle=\|f_+\|^2-\|f_-\|^2.
\end{equation}
Zero expectation means equal weights in these two orthogonal sectors.
It is not a fixed-point condition: a $J$-fixed unit vector has expectation
one. This difference is decisive when a symmetry axis or a critical
line is proposed as an explanation of mass-root balance.

\subsection{Realification and the meaning of an adjoint}
Suppose instead that an inherited $J$ is antiunitary on a complex
Hilbert space. It obeys $J(if)=-iJf$, so $(I+J)/2$ is generally not
complex-linear. Regard the same vectors as a real Hilbert space with
inner product $\operatorname{Re}\langle f,g\rangle$. An antiunitary
involution then becomes an orthogonal real-linear involution, and the
projection calculation is valid on that real space. What changes is
the allowed class of operators and subspaces.

One can also complexify a real orthogonal realization and extend its
involution complex-linearly. This produces a different comparison
construction, in which multiplication by $i$ commutes with $J$.
These alternatives cannot be combined in an argument by keeping only
the symbol $J$. Later formulas such as $-iJA$ require the stated
complex-linear extension and a selfadjoint anticommuting $A$ on the
same positive Hilbert space. They are not formulas for an unspecified
antiunitary operation.

The same discipline applies to tensors. A spacetime metric contracts
spacetime indices; an internal Hermitian metric compares internal
amplitudes; a generation metric compares generation coordinates.
Their signatures, dimensions and transformation rules differ. Even
when their numerical matrices happen to coincide in one chart, an
identification must explain how the corresponding spaces and adjoints
are transported. The QAM context records discussed next serve to make
these choices visible at the point where they enter an inference.

\subsection{The physical interface as an additional piece of data}
A physical readout consists of a domain of admitted states together
with a map to observables. For the mass branch this includes a
generation assignment, positivity of principal roots and a dimensional
normalization. For a geometric branch it includes an invertible
coframe or a metric, its differentiation rules and its coupling to
the equations of motion. These requirements are stronger than the
existence of a vector-space isomorphism.

The common-carrier proposal can therefore be investigated in two
directions. Forward construction starts from a represented carrier
and computes its readout. Backward analysis starts from a required
property of the observable and asks which carrier operations could
produce it. The mass-root conditions and the projected-autonomy test
in this paper are backward constraints of this second kind. They do
not discard the common-carrier hypothesis; they specify the extra
mathematical work needed to give it physical content.
\section{\raggedright Statement language and reproducibility}
The Companion uses QAM as an object and assertion language, MML for
typed instructions and MMA for their execution semantics. For the
physical arguments in this paper, their essential role is to keep an
assertion attached to its domain, assumptions, calculation and remaining
interpretation requirements. A guard denotes a condition still needed
for a stated inference. A finite calculation and a general proof have
different scopes, even when both concern the same formula.

The projection identity provides a complete elementary example. Its
inputs are a Hilbert space, $J^2=I$, $J^*=J$ and the definition of $C$;
the calculation in the preceding section proves $C^2=C$. An exact
finite instance is
\begin{equation}
 P=\frac12\begin{pmatrix}1&1\\1&1\end{pmatrix},\quad
 M=\frac14\begin{pmatrix}5&3\\3&5\end{pmatrix},\quad
 N=\begin{pmatrix}1&0\\0&2\end{pmatrix}.
\end{equation}
Direct multiplication gives $P^2=P$, $[M,P]=0$, and
\begin{equation}
 [N,P]=\frac12\begin{pmatrix}0&-1\\1&0\end{pmatrix}\ne0.
\end{equation}
An exact rational check therefore confirms the projection and one
commuting transport while rejecting the other transport as a commuting
example. This rejection is a control result, not a failure of the
projection theorem. None of these matrix equations identifies a native
field, a physical transport or a mass operator.

Ordinary mathematical notation is retained throughout the main text.
The accompanying assertion records specify precise source versions,
prerequisites and finite checks; they allow the chain of reasoning to
be inspected without treating content identity as proof. The original
mass-readout arguments are collected in Appendix~A, while the separate
audit records retain the detailed provenance and unresolved obligations.
\subsection{An assertion is more than a displayed formula}
Consider the familiar expression $C^2=C$. The formula is ambiguous
until the domain of $C$, the scalar field, the composition operation
and the assumptions on its defining involution are fixed. It also
does not state whether the assertion is adopted, derived generally,
or checked only for a particular matrix. A physical application adds
the question of what is measured by the range of $C$.

For exposition, write an assertion record as
\begin{equation}
 r=(\varphi,\mathcal F,s,W,G),
\end{equation}
where $\varphi$ is the statement, $\mathcal F$ its declared context,
$s$ its evidence status, $W$ its sources and premises, and $G$ its
remaining guards. This tuple summarizes the record's role here; it
is not an additional literal MML instruction. Two records with the
same displayed formula can describe different assertions when their
contexts or premises differ.

QAM is the project's language for such objects and assertions. MML
specifies operations on the declared objects. MMA interprets the
supported instructions and returns execution information. The
implementation used in the examples is a finite interface, not an
unrestricted interpreter for every symbolic statement in the Companion.
Consequently an indexed formula is not automatically executable, and
a successful execution is not automatically a proof of its physical
interpretation.

An object can carry explicit index-role data. The implemented interface
records an index label, its side, variance, role and declared space.
The side and role must be a permitted pair in that interface. A
contraction must also specify which spaces are paired. For instance,
two upper structural indices do not define a scalar product until
a compatible comparison tensor is supplied. Merely printing repeated
labels cannot replace that tensor. This is especially relevant when
an internal comparison is proposed as a spacetime contraction.

\subsection{Derivations, combinations and frame changes}
Three simple operations illustrate the distinction between manipulating
records and proving their contents. A derivation step attaches a
conclusion to its premise record. A join combines records in a common
context. A conversion transports a record through an explicitly
declared map. In the independent provenance checker these appear as
the constructors \texttt{step}, \texttt{join} and \texttt{conv}.

The checker uses an ordered status vocabulary: open, reading,
conditional, computed and proved. This is a protocol ordering used
to prohibit unsupported promotion. It is not a universal ranking of
scientific confidence. In particular, a conditional general theorem
and an exact computation of one example answer different questions.
A valid record tree cannot obtain a stronger status simply by
combining or converting its premises. Its witness inputs remain
tracked, and a frame conversion must name a declared comparison map.

These structural properties have formal record-level statements in
the project's Lean implementation. Their scope is important. A
theorem that well-formed records do not lose witness identifiers
does not prove that a witness describes the correct physical system.
Similarly, a declared map can be type-correct before its metric or
observable-preservation properties have been proved. The semantic
conditions must appear as their own premises or guards.

In the LaTeX interface, a frame declaration introduces an identifier,
a map declaration specifies its source and target, and a formula
declaration supplies a formula, context and status. The actual macro
names include \texttt{qamframe}, \texttt{qammap},
\texttt{qamformula} and \texttt{qamphysconvert}. These are document
commands rather than assertions that a physical conversion exists.
They allow the same ordinary equation to be displayed with its
explicit profile when taken outside its local context. A laboratory
formula and a formula in a boosted-observer profile must therefore
retain the information identifying their different preparation and
comparison data.

\subsection{A complete physical example: information lost by a readout}
Take classical electromagnetic fields in units $c=1$ and form the
three-by-two matrix $F=[\boldsymbol E\mid\boldsymbol B]$. On nonzero
fields define the normalized Gram readout
\begin{equation}
 \rho(F)=\frac{F^{\mathsf T}F}{\operatorname{Tr}(F^{\mathsf T}F)}.
\end{equation}
This positive two-by-two matrix records relative norms and the scalar
product of the two vectors. It discards their common scale and their
orientation in physical space. These omissions matter under an
observer transformation whose direction is specified in the laboratory.

For a boost of velocity $v$ in the positive $z$ direction, consider
the pair of preparations
\begin{equation}
 \boldsymbol E=2e_x,\qquad \boldsymbol B=b e_y,\quad b=\pm1.
\end{equation}
Both have the same input register $\rho=\operatorname{diag}(4/5,1/5)$.
The ordinary field transformation gives
\begin{equation}
 \boldsymbol E'_b=\gamma(2-vb)e_x,\qquad
 \boldsymbol B'_b=\gamma(b-2v)e_y,
 \quad \gamma=(1-v^2)^{-1/2}.
\end{equation}
At $v=1/4$, the common factor cancels from the normalized Gram
readout and the two outputs are
\begin{equation}
 \rho(F'_{+1})=\operatorname{diag}(49,4)/53,\qquad
 \rho(F'_{-1})=\operatorname{diag}(9,4)/13.
\end{equation}
Thus there is no single-valued boosted readout depending only on
the original normalized register for this family. A map on that
register would need to assign two different outputs to one input.
This is an elementary failure of information sufficiency, not a
failure of the Lorentz transformation on the full fields.

The example explains why a frame label alone cannot repair a reduced
description. One must either retain additional preparation data or
replace a single predicted value by the set of outputs compatible
with the data actually available. This same distinction will later
appear in the hidden-lift term of projected dynamics.

\subsection{An exact interval prediction from an additional measurement}
Allow the laboratory preparation to vary as
\begin{equation}
 \boldsymbol E=2e_x,\qquad
 \boldsymbol B=t e_y+u e_z,\qquad t^2+u^2=1.
\end{equation}
All these fields again have the same input Gram register. Suppose
an extra measurement constrains the signed laboratory component
$t=\boldsymbol B\cdot e_y$ to $[a,b]\subseteq[-1,1]$. At the
same boost speed $v=1/4$,
\begin{equation}
 \boldsymbol E'=\gamma(2-t/4)e_x,\qquad
 \boldsymbol B'=\gamma(t-1/2)e_y+u e_z.
\end{equation}
The vectors remain orthogonal. Divide their squared norms by
$\gamma^2$ and set $A(t)=(2-t/4)^2$. Then
\begin{equation}
 \frac{|\boldsymbol B'|^2}{\gamma^2}
 =(t-1/2)^2+(15/16)(1-t^2)=A(t)-45/16.
\end{equation}
Consequently the first diagonal entry is
\begin{equation}
 p(t)=\frac{A(t)}{2A(t)-45/16}
     =\frac{(8-t)^2}{2(8-t)^2-45}.
\end{equation}
The denominator is strictly positive on the admitted interval.
Moreover $A$ decreases with $t$ and
$dp/dA=-(45/16)/(2A-45/16)^2<0$, so $p$ increases. The exact
image of the measurement interval is therefore
\begin{equation}
 \rho'\in\{\operatorname{diag}(p,1-p):p(a)\leq p\leq p(b)\}.
\end{equation}
Every intermediate value is attained by some preparation in the
specified family. This is a deterministic interval statement; no
probability distribution on $t$ has been assumed.

For $[a,b]=[-1/4,1/4]$ the exact result is
\begin{equation}
 p\in[121/162,\ 961/1202]
      \subset[9/13,\ 49/53].
\end{equation}
The wider interval is the output when only $t\in[-1,1]$ is known.
The additional measurement has therefore narrowed the prediction
without reconstructing an arbitrary orientation. The numbers in
this example are a controlled classical preparation, not fitted
lepton masses or an empirical confirmation of the common-carrier
hypothesis.

\subsection{From the assertion to the MML request and the MMA receipt}
The QAM record for this example declares a map
$\mathrm{Lab}\to\mathrm{Boosted}$ and two conditional assertions:
the preparation, and the interval rule derived from it. In the
implemented request the instruction is \texttt{READOUT}. Its
arguments name the recipe \texttt{em.gram-boost-range}, the preparation
\texttt{E2x-B-unit-yz}, source frame \texttt{Lab}, target frame
\texttt{Boosted}, boost \texttt{z:1/4}, and endpoints
\texttt{-1/4} and \texttt{1/4}. The endpoints are rational strings,
so the adapter can evaluate them without a floating rounding step.

The stored receipt for request \texttt{gram\_boost\_interval} returns
the exact target interval $[121/162,961/1202]$. Its recorded premise
chain links the interval rule to the stated laboratory preparation.
This receipt was retrieved by its content hash and checked through
the read-only evidence interface for this edition. The complete
identifier and source fingerprints belong in the accompanying evidence
files, where a reader can distinguish this saved calculation from
a newly executed or replayed one.

Its execution status is \texttt{continue}, with no outstanding
adapter obligations and no missing recorded premises. The premises
themselves remain conditional. In particular, the same receipt records
that closure is not proved and that no general mathematical relation
is certified by the interpreter. The analytic derivation above supplies
the general interval argument; the receipt verifies the specified
instance and its recorded execution context.

The other execution outcomes have equally limited meanings.
\texttt{guard} identifies an unmet prerequisite for the requested
operation, whereas \texttt{reject} identifies a request that cannot
be accepted under the interpreter's contract. These outcomes are not
statistical likelihoods or statements about the truth of an entire
physical theory. An empty list of additional request guards does not
erase the assumptions in the preparation or the rule.

\subsection{What is preserved when the document is regenerated}
The document store retains immutable assertion versions, source
identities and ordered selections. A reader edition can move an
original proof to an appendix while keeping its exact source
association. A change of prose order does not create a new inference,
and a content hash does not settle whether an imported theorem heading
is justified. The expanded physical argument must continue to state
its mathematical assumptions in the body, rather than relying on a
reader to discover them in an audit log.

This division of work is useful in a multi-branch programme. The main
paper explains a restricted set of connected results. Its evidence
package supplies the reproducible records for those results. The long
Companion retains alternative routes and their history, including
failed or qualified claims. A future change to a preparation or a
comparison tensor must be traced through dependent results before the
corresponding physical conclusion can be reused.
\section{\raggedright Particle and generation assignments}
The algebraic carrier suggests several representation-theoretic
ingredients: unit quaternions provide an $SU(2)$ action, complex scalars
provide phase transformations, and the octonionic setting motivates
colour representations. These ingredients are useful starting points,
but they do not uniquely fix the physical gauge theory. In particular,
a common complex phase is not by itself the electromagnetic charge
operator or a derivation of hypercharge.

The distinction between colour and generation is equally important.
A three-dimensional subspace carrying a colour action cannot be
identified with three generations merely by counting its basis vectors.
The Companion's inherited slot assignments require additional choices
of embeddings, charge operators and the relation between internal
representations and physical states. We therefore retain the gauge and
generation discussion as a conditional assignment problem, not as a
derivation of the complete Standard Model.

A more explicit comparison construction starts from a seed isometry
$T_0:\mathbb C^3\to\mathcal K$ and an orthogonal closure projector
$P_c$ on a declared common carrier. If $T_0^*P_cT_0$ is positive
definite, then
\begin{equation}
 S=P_cT_0(T_0^*P_cT_0)^{-1/2},\qquad S^*S=I_3.
\end{equation}
This constructs three orthonormal closed states under its hypotheses.
It does not determine the seed embedding or explain why those states
are lepton generations. We distinguish $P_c$ from the HC projector $C$:
closure within a selected space and even/odd HC grading are separate
structures. This distinction becomes necessary when a coupling is
used to define mass readout.
\subsection{What a representation construction must specify}
The algebraic question is not whether symbols resembling the three
gauge factors can be placed in a diagram. It is which operators act
on which states, how the actions commute, and which observables are
assigned to their invariant subspaces. A physical gauge interpretation
also requires local transformation laws, connections and the appropriate
matter representations. The compact groups encountered in a carrier
are only part of these data.

For the quaternion factor, write the Hermitian generators as
$T_k=\sigma_k/2$. The Pauli multiplication law gives
\begin{equation}
 [T_j,T_k]=i\epsilon_{jkl}T_l.
\end{equation}
This exhibits the Lie algebra of the unit-quaternion $SU(2)$ action.
It does not decide whether a particular subspace represents a weak
doublet, a spatial spinor or another two-component system. That
decision depends on the remaining representation and coupling data.
Likewise the commuting scalar phase $e^{i\theta}$ acts uniformly on
the carrier until a nontrivial charge operator is specified.

The octonionic ingredient is equally sensitive to the chosen complex
structure. Fix an imaginary unit $e_7$ and let $I_7=L_{e_7}$ denote
left multiplication on the real octonion space. Here $I_7^2=-I$,
so it is a complex structure, not the involution $J$ used for the HC
grading. Relative to it, the real eight-dimensional space can be
read as a complex four-dimensional space. The familiar singlet/triplet
organization associated with the stabilizer of $e_7$ motivates the
colour comparison used in the Companion. All identifications in this
paragraph refer to that selected internal complex structure.

On the complexification, the corresponding complex-linear projections
are
\begin{equation}
 \Pi_\pm=\frac{I\mp iI_7}{2},\qquad
 \Pi_\pm I_7=\pm i\Pi_\pm.
\end{equation}
Complex coefficient conjugation exchanges these projections. They
must not be substituted for $C$ or for a closure projector without
a map that identifies the relevant domains and meanings. The same
letter ``projector'' can otherwise hide three different decompositions:
internal complex structure, HC parity and physical closure.

\subsection{A conditional ladder construction and its exact counting}
The charge construction retained in the Companion can be separated
into an algebraic hypothesis and the consequences that follow from it.
In its specified octonion multiplication convention, introduce
\begin{equation}
 \alpha_a=\tfrac12(L_{e_a}+iL_{e_7e_a}),\quad a=4,5,6,
 \qquad N=\sum_{a=4}^6\alpha_a^*\alpha_a,
 \qquad Q_{\rm ch}=N/3.
\end{equation}
The represented Clifford multiplication and adjoint must give the
canonical anticommutation relations
\begin{equation}
 \alpha_a\alpha_b+\alpha_b\alpha_a=0,\qquad
 \alpha_a\alpha_b^*+\alpha_b^*\alpha_a=\delta_{ab}I.
\end{equation}
Once these relations and a cyclic vacuum annihilated by the three
$\alpha_a$ are supplied, the occupation counting is elementary.
Each creation operator can occur at most once, because its square
vanishes. The states with $k$ occupied modes therefore have dimension
$\binom3k$, and $N$ has eigenvalue $k$ on that subspace. The charge
comparison consequently has the spectrum
\begin{equation}
 \begin{array}{c|rrrr}
 k&0&1&2&3\\\hline
 \text{multiplicity}&1&3&3&1\\
 Q_{\rm ch}&0&1/3&2/3&1
 \end{array}
\end{equation}
The total dimension is eight, the dimension of the three-mode fermionic
Fock representation. These multiplicities are consequences of the
chosen ladder algebra. They do not establish three lepton families.

An internal $SU(3)$ action that rotates the creation operators and
preserves the vacuum acts on the one-occupation space as a triplet,
and on the two-occupation space as its conjugate, since these are
the first and second exterior powers of a three-dimensional complex
space. Vacuum and full occupation are singlets. This gives a precise
representation question: is the represented colour action on the
proposed octonionic carrier the action required by these operators?
It must be checked with the chosen multiplication and adjoint; the
counting table alone is not that check.

The Companion's historical calculations motivate this compatibility,
but the current paper does not turn a preserved numerical calculation
into a new formal theorem. Its physical use still needs the embedding
of particles, their chirality and the electroweak charge assignments.
In particular a number operator with the displayed eigenvalues does
not, by itself, derive hypercharge or the full weak representation.

This decomposition also explains a correction to the literal slot
dictionary. With the inner complex structure, vectors such as
$e_4,e_5,e_6$ can form a complex triplet basis only because their
$I_7$-images supply the corresponding real partners. Counting the
images as another independent physical triplet double-counts the
same complex degrees of freedom. Changing to the commuting outer
complex unit changes the decomposition into charge eigenvectors.
State labels must follow that change; an unqualified equality of
charges across the proposed frames is not justified.

\subsection{The independent generation space}
Three colours within one representation and three generations of a
given particle are different multiplicities. If the colour action is
nontrivial on a three-dimensional subspace, that action does not
automatically act as a generation permutation. A generation assignment
needs a second declared structure and its relation to the carrier.
The closure-seed construction supplies a conditional way to express
that requirement without identifying raw octonion coordinates with
mass eigenstates.

Let $T_0:\mathbb C^3\to\mathcal K$ be an isometric seed embedding,
and let $P_c$ be an orthogonal projector. Define
$G_0=T_0^*P_cT_0$. If $G_0>0$, then the states in
\begin{equation}
 S=P_cT_0G_0^{-1/2}
\end{equation}
are closed and orthonormal:
\begin{equation}
 P_cS=S,\qquad
 S^*S=G_0^{-1/2}T_0^*P_c^2T_0G_0^{-1/2}=I_3.
\end{equation}
If $G_0$ has rank below three, projection has removed an independent
seed direction and cannot produce three independent states by this
formula. If $G_0$ is not diagonal, normalization already mixes the
original seed slots. Thus even a successful construction does not
identify each initial slot separately as a physical mass eigenstate.

Suppose also that $[P_c,J]=0$ and $JT_0=-T_0$. The same calculation
gives $JS=-S$. This realizes three closed odd inputs. A nonzero
operator restricted to this odd space preserves oddness and commutes
with the restricted $J=-I$. It cannot also anticommute with $J$.
An HC-odd coupling must connect opposite parity sectors; an operator
that merely has odd outputs is a different object.

The distinction supplies a useful diagnostic before computing any
spectrum. A proposed Dirac or mass coupling should specify whether
it acts within a sector, changes HC grading, changes closure sector,
or combines these operations. Its matrix blocks and domains determine
the answer. Calling the operator ``odd'' or giving it a familiar
letter does not determine which of these statements is intended.

\subsection{A represented transition and the choice of mass basis}
Take a bounded dimensionless comparison pacer $\widetilde R$ on the
same carrier and define
\begin{equation}
 D=(I-P_c)\widetilde R P_c,\qquad B=DS,\qquad Q=B^*B.
\end{equation}
Then $Q\geq0$. If a mass law $M_{\rm mass}=\kappa Q$ is independently
adopted, its candidate eigenmasses are $\kappa q_i$, where $q_i$
are the eigenvalues of $Q$. A fixed-slot readout instead uses
$\kappa Q_{ii}=\kappa\|B\varepsilon_i\|^2$. These prescriptions
need not coincide.

For example, with $\boldsymbol1=(1,1,1)^{\mathsf T}$, let
\begin{equation}
 Q=\tfrac13\boldsymbol1\boldsymbol1^*.
\end{equation}
Its three diagonal entries equal $1/3$, but its spectrum is $(1,0,0)$.
The fixed-slot mass-root ratio would be $1/3$, whereas the spectral
ratio would be $1$. These are different observables of one positive
matrix. The example excludes the inference from equal diagonal
couplings to degenerate physical eigenmasses.

There is also a rank requirement. Since
$\langle v,Qv\rangle=\|Bv\|^2$, the kernels of $Q$ and $B$
coincide, and
\begin{equation}
 \operatorname{rank}Q=\operatorname{rank}B.
\end{equation}
Three strictly positive candidate masses require three independent
transition outputs. A single scalar output channel has rank at most
one, regardless of how many input labels it is given. An operator-valued
channel may have higher dimension, but that dimension must be represented
rather than inferred from the name of a probe channel.

\subsection{What mixing angles would require}
Mixing is a relative-basis observable. To make the requirement explicit,
consider two Hermitian generation comparisons $Q_1,Q_2$ with specified
physical meanings. Let $V_1,V_2$ diagonalize them in the same generation
space. Their relative unitary matrix is
\begin{equation}
 U_{12}=V_1^*V_2.
\end{equation}
Changing the common generation coordinates by $F$ replaces the two
diagonalizers by $FV_1,FV_2$ and leaves $U_{12}$ unchanged. Rephasing,
ordering and eigenvalue degeneracies require their usual separate
conventions before individual parameters can be interpreted.

This construction identifies the data a mixing prediction would need:
two independently represented operators, their common kinetic metric,
their assignment to physical species, and a rule fixing their relative
eigenspaces. A common cyclic action or the dimension three does not
provide those data. In a degenerate eigenspace the basis is free, so
apparent mixing angles obtained from an arbitrary diagonalizer are
not predictions. The mass-readout analysis below therefore treats
the spectrum, its physical scale and its generation basis as distinct
questions, rather than allowing one successful matrix pattern to
stand in for all three.
\section{\raggedright Mass readout and consistency constraints}
\paragraph{A necessary norm relation.}
Let $x_i=\sqrt{m_i}\geq0$, with $x\ne0$, and define
\begin{equation}
 K=\frac{\sum_i x_i^2}{(\sum_i x_i)^2},\qquad
 A=\frac{x_0+x_1+x_2}{3},\qquad x=A(1,1,1)+d.
\end{equation}
Since $\sum_i d_i=0$, the constant part and $d$ are orthogonal, so
\begin{equation}
 K=\frac13+\frac{\|d\|^2}{9A^2},\qquad
 K=\frac23\ \Longleftrightarrow\ \|d\|^2=3A^2.
\end{equation}
Thus the Koide value requires equality of the symmetric and traceless
squared norms. Cyclic symmetry preserves these norms separately and
does not fix their ratio. For example, $x=(1,2,3)$ gives $K=7/18$,
while equal roots give $K=1/3$.

\paragraph{A conditional HC representation.}
On Euclidean three-root space set $e=(1,1,1)/\sqrt3$ and let $J_m$
be an orthogonal involution. For $C_m=(I+J_m)/2$, suppose
$x=\alpha C_me$ is nonzero and componentwise nonnegative, with
$\alpha>0$. Writing $h=\langle e,J_me\rangle$ gives
\begin{equation}
 \|C_me\|^2=\frac{1+h}{2},\qquad
 K=\frac{2}{3(1+h)}.
\end{equation}
Consequently $h=0$ is the additional condition for Koide. The projector
axioms permit $J_m=I$, for which $h=1$ and $K=1/3$; they do not select
the desired value. Choosing $J_m$ after specifying the target roots
cannot supply the missing physical selection rule.

An explicit route from the original carrier to this representation
would require a bounded coisometry $B:H\to\mathbb R^3$, $BB^*=I_3$,
whose kernel is $J$-invariant. Then $J_m=BJB^*$ is an orthogonal
involution and $BC=C_mB$. If an independently supplied seed satisfies
$Bf_0=e$ and the physical root law is $x=\alpha BCf_0$, the criterion
becomes
\begin{equation}
 \langle B^*e,JB^*e\rangle=0.
\end{equation}
Neither the name of a readout nor the number of its components supplies
$B$, its metric control, this seed law or the orthogonality selection.
Moreover, a purely odd input obeys $Cf=0$ and cannot yield nonzero roots
through this direct compensated linear route.

\paragraph{Positive nonlinear readout is a separate requirement.}
For two complex amplitudes $\chi=(u,v)$ define
$X_0=|u|^2+|v|^2$, $X_1=|u|^2-|v|^2$,
$X_2=2\operatorname{Re}(\bar uv)$ and
$X_3=2\operatorname{Im}(\bar uv)$. Then
$X_0^2=X_1^2+X_2^2+X_3^2$. With the three-cycle $U$, an illustrative
Hermitian root operator is
\begin{equation}
 L_\eta=X_0\frac{I_3}{\sqrt3}
 +\eta X_1\frac{U+U^*}{\sqrt6}
 +\eta X_2\frac{i(U-U^*)}{\sqrt6},\qquad \eta\geq0.
\end{equation}
Only where $L_\eta\geq0$ may its eigenvalues be read as principal
roots. There
\begin{equation}
 K=\frac{1+\eta^2[1-(X_3/X_0)^2]}3.
\end{equation}
Positivity on every input requires $\eta\leq1/\sqrt2$ and gives
$K\leq1/2$. The isometric coefficient $\eta=1$ can reach $2/3$
on restricted positive phase chambers, but the preparation and
preservation of such a chamber is additional physical information.
Squaring a signed $L_\eta$ does not remove this issue: its physical
principal root becomes $|L_\eta|$, changing the root sum.

The recent HC-balance comparison illustrates the transfer problem
more directly. If the two-mode grading is $J=\operatorname{diag}(1,-1)$,
the balanced state $(1,i)/\sqrt2$ gives $X_1=X_2=0$ and hence $K=1/3$
for this readout. More generally, an antipodal pair of root operators
$cI_3\pm D$, with $c>0$ and $\operatorname{Tr}D=0$, cannot both
be positive Koide readouts. Both positive implies eigenvalues of $D$
in $[-c,c]$ and $\operatorname{Tr}D^2\leq2c^2$, so both ratios
are at most $5/9$. This is a restriction on this isotropic readout
pair, not on arbitrary mass mechanisms. It shows why state balance
and the positive three-generation transfer must be derived together.
\subsection{Why a transition norm is a new readout law}
For a closed seed $P_c\psi=\psi$, the closure commutator satisfies
\begin{equation}
 [\widetilde R,P_c]\psi=(I-P_c)\widetilde R\psi=D\psi.
\end{equation}
The right-hand side lies in the orthogonal complementary space. Hence
$\langle\psi,[\widetilde R,P_c]\psi\rangle=0$, even if its
norm is nonzero. Replacing a commutator expectation by
$\|D\psi\|^2$ therefore changes the observable. It is a legitimate
candidate only when that change is stated and physically justified.

The positive matrix $Q=S^*D^*DS$ measures a tendency to leave the
selected closure space under the supplied operator. It does not
assert that an evolving physical state actually leaves that space.
The operator might be a diagnostic rather than the time generator.
The physical mass law and the preservation of the admitted state
domain must be supplied independently.

The distinction between closure and HC parity gives another necessary
test. Under $[P_c,J]=0$, define
\begin{equation}
 R_o=\tfrac12(\widetilde R-J\widetilde RJ),\qquad
 D_o=\tfrac12(D-JDJ)=(I-P_c)R_oP_c.
\end{equation}
Only $D_o$ changes HC grading. A pacer commuting with $J$ has
$R_o=0$ and cannot supply this channel, even if it has nonzero
closure defect. A selfadjoint comparison coupling can be formed as
$A_o=D_o+D_o^*$, but forming it is an additional operation; its
appearance in the native equations is not established by its definition.

\subsection{Closure penalties do not determine the missing channel}
Let $M$ be a bounded represented matrix and split the carrier into
the range and kernel of $P_c$. Write
\begin{equation}
 M=\begin{pmatrix}A&F\\G&H\end{pmatrix},\quad
 K_c=[M,P_c]=\begin{pmatrix}0&-F\\G&0\end{pmatrix},\quad
 E_c=\begin{pmatrix}0&F\\G&0\end{pmatrix}.
\end{equation}
Direct block multiplication gives
\begin{equation}
 K_c^*K_c=E_c^*E_c=
 \begin{pmatrix}G^*G&0\\0&F^*F\end{pmatrix}.
\end{equation}
Thus a minimal positive penalty
$H_c=\rho K_c^*K_c+\sigma E_c^*E_c$ is block diagonal and
commutes with $P_c$. Its compression can measure the norm of a
transition generated by $M$, but the penalty itself has no off-diagonal
closure channel. Positivity of the penalty cannot replace the missing
represented pacer block.

If the pacer were a norm limit of polynomials in an exactly closed
$M,M^*$, then it would commute with $P_c$ as well and its transition
would vanish. Conversely, exact closure of $M$ says nothing about
an independently specified pacer. One needs a bridge between those
operators before transferring a closure statement from one to the other.

There is a particularly sharp freedom result even when the pacer is
unitary. Given any positive contraction $0\leq Q\leq I_3$, define
\begin{equation}
 R_Q=\begin{pmatrix}\sqrt{I_3-Q}&-\sqrt Q\\
                         \sqrt Q&\sqrt{I_3-Q}\end{pmatrix},\qquad
 P_c=\operatorname{diag}(I_3,0).
\end{equation}
The square roots commute because both are functions of $Q$.
Block multiplication gives $R_Q^*R_Q=I_6$, and the transition from
the first three coordinates to the complement has Gram matrix $Q$.
Therefore unitarity, positive readout and the dimension three leave
every positive contraction spectrum available. These assumptions
do not select a hierarchy or the Koide radius.

\subsection{When a Gram law follows from a gapped block}
A possible physical route can be expressed more concretely at a
frozen background. Assume a canonically normalized Hermitian mass
operator on a retained three-dimensional space and its complement:
\begin{equation}
 \Omega_{\rm tr}=\begin{pmatrix}A&V^*\\V&H_\perp\end{pmatrix},
 \qquad A=A^*,\quad H_\perp=H_\perp^*\geq\Delta I,
 \quad \Delta>0.
\end{equation}
For an eigenvalue $|\lambda|<\Delta$, solve the lower equation to obtain
\begin{equation}
 y=-(H_\perp-\lambda I)^{-1}Vx,
\end{equation}
and substitute into the upper equation:
\begin{equation}
 [A-V^*(H_\perp-\lambda I)^{-1}V]x=\lambda x.
\end{equation}
This is an energy-dependent Schur complement. The resolvent identity
controls the error made by replacing it by its static value:
\begin{equation}
 \|A_{\rm eff}(\lambda)-A_{\rm eff}(0)\|
 \leq\frac{\|V\|^2|\lambda|}{\Delta(\Delta-|\lambda|)}.
\end{equation}
The bound uses the two inverse norms
$\|(H_\perp-\lambda I)^{-1}\|\leq(\Delta-|\lambda|)^{-1}$
and $\|H_\perp^{-1}\|\leq\Delta^{-1}$. It is a controlled
comparison only in the stated gap regime.

With the further assumptions $A=0$, $H_\perp=\Delta I$ and
$V=\mu B$, each eigenvalue $q_i$ of $Q=B^*B$ gives a light
signed eigenvalue $-m_i$, where
\begin{equation}
 m_i=\frac{\sqrt{\Delta^2+4\mu^2q_i}-\Delta}{2}
 =\frac{\mu^2}{\Delta}q_i
  +O\left(\frac{\mu^4q_i^2}{\Delta^3}\right).
\end{equation}
Thus a Gram mass-magnitude law can be a leading approximation in
a canonical Hermitian Dirac-type reading. The signed effective
operator is negative in this comparison; its interpretation as a
positive physical mass magnitude is an additional reading, not
a sign that can be dropped without comment.

The exact eigenstate is dressed. Its complementary norm satisfies
\begin{equation}
 \frac{\|y\|^2}{\|x\|^2}=\frac{m_i}{\Delta+m_i},\qquad
 \|y\|^2=\frac{m_i}{\Delta+2m_i}
\end{equation}
after normalization. Enforcing strict closure $y=0$ would instead
require $Vx=0$. A nonzero light mass produced in this zero-direct-block
mechanism is consequently incompatible with exact closure in the
original retained projector. Approximate closure or a separately
defined dressed spectral projector are possible alternatives, but
they are different state constraints.

\subsection{Cyclic compatibility fixes a form, not its coefficients}
Suppose a unitary carrier operation $W$ and the generation three-cycle
$U$ satisfy
\begin{equation}
 WS=SU,\qquad WP_c=P_cW,\qquad W\widetilde R=\widetilde RW.
\end{equation}
Then $BU=WB$ and hence $QU=UQ$. The positive root $L=\sqrt Q$
also commutes with $U$. In three dimensions the three distinct Fourier
eigenvalues of $U$ imply that every such Hermitian root has the form
\begin{equation}
 L=aI_3+zU+\overline zU^*,\qquad a>0,\quad z=re^{i\theta}.
\end{equation}
Its eigenvalues are
\begin{equation}
 \ell_k=a+2r\cos(\theta+2\pi k/3).
\end{equation}
Trace orthogonality gives
\begin{equation}
 \operatorname{Tr}L=3a,\qquad
 \operatorname{Tr}L^2=3a^2+6r^2,\qquad
 K=\frac13+\frac{2r^2}{3a^2}.
\end{equation}
The Koide value therefore requires $r=a/\sqrt2$. Cyclic compatibility
does not supply this coefficient ratio. Conversely, once the native
$Q$ were represented, its coefficients would be determined without a
mass fit by
\begin{equation}
 a=\tfrac13\operatorname{Tr}\sqrt Q,\qquad
 z=\tfrac13\operatorname{Tr}(U^*\sqrt Q).
\end{equation}
The native burden is thus a joint operator calculation, not the
selection of a familiar trigonometric parametrization.

At the Koide radius positivity gives three phase chambers,
\begin{equation}
 \theta\in\bigcup_{j=0}^2
 [2\pi j/3-\pi/12,\ 2\pi j/3+\pi/12]
 \pmod{2\pi}.
\end{equation}
They follow by intersecting the inequalities
$\cos(\theta+2\pi k/3)\geq-1/\sqrt2$ for all three $k$.
At the boundary $\theta_b=3\pi/4$, write
$\theta=\theta_b-\delta$ inside the chamber. The small root is
$\ell_0=a\delta+O(\delta^2)$, so its candidate mass is
$m_0=\kappa a^2\delta^2+O(\delta^3)$. A light eigenmass can
therefore arise near a boundary in this comparison. Neither the
distance to the boundary nor its dynamical preservation has yet been
selected by the native theory.

\subsection{Deriving the positivity obstruction in the quadratic model}
For $\eta\geq0$, the quadratic root model can be written as a linear function of the
two-mode density matrix $\rho=\chi\chi^*$. In the Fourier basis of
$U$, its signed root eigenvalues are
\begin{equation}
 \ell_k=\frac{X_0}{\sqrt3}
 +\sqrt{\frac23}\eta[X_1\cos\varphi_k-X_2\sin\varphi_k],
 \qquad \varphi_k=2\pi k/3.
\end{equation}
Equivalently $\ell_k=\sqrt3\operatorname{Tr}(F_k\rho)$, where
\begin{equation}
 F_k=\frac{I_2+\sqrt2\eta(\cos\varphi_k\sigma_z-
                               \sin\varphi_k\sigma_x)}{3}.
\end{equation}
The three $F_k$ sum to $I_2$. The Pauli combination has eigenvalues
$\pm1$, so $F_k$ has eigenvalues $(1\pm\sqrt2\eta)/3$.
Positivity for every $\rho\geq0$ is therefore equivalent to
$\eta\leq1/\sqrt2$. Necessity follows by testing a negative
eigenvector if this condition fails; sufficiency follows because
the trace of a product of two positive matrices is nonnegative.
This is the precise sense in which a globally positive linear-density
readout in this family cannot achieve the intended isometric radius.

Replacing the input metric by the pullback metric of this map would
make a chosen nonzero $\eta$ an isometry by definition. It would
not derive its physical value. Nor does covariance under the three-cycle
fix the relative singlet and doublet weights: the constant direction
and the zero-sum plane are separate invariant pieces. A comparison
tensor can assign them different positive weights while preserving
the cyclic symmetry.

The antipodal $5/9$ bound is similarly a property of the actual
root eigenvalues. If $cI_3\pm D$ are both positive, the three
eigenvalues $d_i$ obey $|d_i|\leq c$ and $\sum_i d_i=0$.
The maximum of $\sum_i d_i^2$ on this planar cube section is
$2c^2$, attained at permutations of $(c,0,-c)$. Consequently
\begin{equation}
 \frac{\operatorname{Tr}[(cI_3\pm D)^2]}
      {[\operatorname{Tr}(cI_3\pm D)]^2}
 =\frac{3c^2+\operatorname{Tr}D^2}{9c^2}\leq\frac59.
\end{equation}
Selecting one allowed positive branch can evade the pairwise
requirement. A physical explanation then has to supply that branch
selection rather than infer it from the involution alone.

\subsection{Phase overlap: root averaging and mass averaging differ}
The published overlap comparison provides another explicit test of
what is being averaged. In one common generation basis, let
\begin{equation}
 x_{jk}=a[1+\sqrt2\operatorname{Re}(\omega^ke^{i\delta_j})],
 \quad \omega=e^{2\pi i/3},\quad
 w_j\geq0,\quad\sum_jw_j=1.
\end{equation}
Each preparation is assumed to have nonnegative roots and the same
mean $a$. Averaging the roots before squaring gives
\begin{equation}
 \overline x_k=a[1+\sqrt2\operatorname{Re}(\omega^k Z)],
 \qquad Z=\sum_jw_je^{i\delta_j},\qquad
 K=\frac{1+|Z|^2}{3}.
\end{equation}
For two equally weighted phases $\delta_0\pm h$, this reduces the
radius by $\cos h$ when $|h|<\pi/2$, leaving the central phase.
It follows from a vector average on the unit circle, not from a
new dynamical mass equation.

Averaging squared roots gives a different answer:
\begin{equation}
 \sum_jw_jx_{jk}^2=\overline x_k^2+
                      \sum_jw_j(x_{jk}-\overline x_k)^2.
\end{equation}
The variance term cannot be removed by calling both procedures a
superposition. A coherent coupling of two paths is yet another
specified construction. For $B(\phi)=B_1+e^{i\phi}B_2$,
\begin{equation}
 B(\phi)^*B(\phi)=B_1^*B_1+B_2^*B_2+
 e^{i\phi}B_1^*B_2+e^{-i\phi}B_2^*B_1.
\end{equation}
Here relative phases can affect the operator spectrum through cross
terms. A common phase multiplying all of $B$ cancels from its Gram
matrix. Changing only the prepared probe can change spectral weights
without moving the operator poles. These alternatives must be
distinguished before a transport interpretation is assigned.

An inverse reconstruction can find phase mixtures consistent with
specified masses and an assumed intermediate gap. Its feasibility
region is useful for excluding combinations of assumptions, but the
reconstructed phases and gap are not independent predictions. The
native operator, the admitted preparations and the rule for combining
paths must be fixed before the resulting spectrum is compared with
data as a test of the model.
\section{\raggedright Projected dynamics and geometry}
The same interface issue appears in the dynamics. Suppose a finite
realification of the full state satisfies $\dot u=Gu$ and let
$r=T(t)u$ be a surjective real-linear extraction. For a right inverse
$R_s$, write $u=R_sr+z$, with $Tz=0$. Differentiation gives
\begin{equation}
 \dot r=(\dot T+TG)R_sr+(\dot T+TG)z.
\end{equation}
The observable equation is independent of the discarded lift precisely
when
\begin{equation}
 (\dot T+TG)\ker T=0.
\end{equation}
This condition is additional to an algebraic projector identity. In its
absence, two full states with the same readout can have different
readout derivatives. Eliminating the hidden variables may then produce
memory terms rather than an autonomous local equation.

The Companion's geometric programme studies induced coframes, local
field comparisons and compatible constraint data. Finite normal-source
calculations can test whether prescribed tangential data determine a
normal increment. Such a calculation is not yet a regional existence
theorem, a derivation of the native action or an identification of an
internal positive comparison metric with a spacetime metric. The
reconstruction must respect common domains, conservation conditions
and the ordering of the original operators.

This identifies a concrete link between the mass and geometry branches:
both require a represented extraction with controlled discarded degrees
of freedom. A formal master equation supplies an organizing equation;
well-posed evolution and the physical meaning of its projected fields
are further obligations.
\subsection{The formal full-state equation and its actual assumptions}
The retained architecture organizes the probe through an ordered equation
\begin{equation}
 (i\Gamma^\mu D_\mu-\Omega)\Psi=0,\qquad
 D_\mu=\partial_\mu+\mathcal B_\mu.
\end{equation}
Its ingredients include an internal connection, a closure projector,
a mismatch field and proposed triadic couplings. Specifying their
names does not determine their joint matrix elements on the generation
states. It also does not ensure the adjoint relations needed for a
real action or a selfadjoint time evolution.

On a chosen time chart with invertible $\Gamma^0$, the equation can
be rearranged algebraically as
\begin{equation}
 \partial_t\Psi=\mathcal G\Psi,
 \quad \mathcal G=-\mathcal B_0-(\Gamma^0)^{-1}
 \sum_{a=1}^3\Gamma^a(\partial_a+\mathcal B_a)
              -i(\Gamma^0)^{-1}\Omega.
\end{equation}
This expression retains the order of noncommuting factors. Turning it
into an evolution problem requires a domain, boundary or initial data,
regularity and an appropriate adjoint. In a finite modal comparison
these questions reduce to represented matrices; in a field theory they
cannot be removed by matrix notation.

The extraction condition given above tests an additional issue even
when a full evolution is already supplied. The reduced variables must
retain sufficient information to determine their own derivative.
The kernel of the extraction is the exact place where this condition
can fail. This is the dynamical counterpart of the electromagnetic
example, in which a Gram register omitted the orientation needed for
the transformed readout.

\subsection{A hidden-lift counterexample with equal normalization}
The dependence on discarded variables is not only a possibility for
large nonlinear systems. Consider a two-component complex amplitude
$f=x+ir$ evolving by $\dot f=-i\sigma_x f$, with $x,r$ real.
Then
\begin{equation}
 \dot x=\sigma_xr,\qquad\dot r=-\sigma_xx.
\end{equation}
Extract only $r$. Fix $0<b<1$ and compare the initial states
\begin{equation}
 r=(b,0)^{\mathsf T},\qquad
 x_\pm=(\pm\sqrt{1-b^2},0)^{\mathsf T}.
\end{equation}
Both full states have unit norm and the same extracted vector. Their
extracted derivatives are
\begin{equation}
 \dot r_\pm=(0,\mp\sqrt{1-b^2})^{\mathsf T}.
\end{equation}
The polar angle of $r$ therefore has opposite initial angular
velocities in the two cases. Equal extracted data and equal full
norm have not defined an autonomous angular law. This example is an
ordinary finite Hermitian evolution; it does not rely on an unstable
Hamiltonian or a noisy numerical reduction.

The linear criterion pinpoints the failure. If $T$ extracts $r$,
then $TG$ does not vanish on $\ker T$, because the discarded real
part drives $\dot r$. Adding a time-dependent extraction introduces
$\dot T$ rather than eliminating that condition. A moving frame may
be physically appropriate, but its derivative must be part of the
calculation.

\subsection{Exact elimination retains memory and initial-state data}
For a constant Hermitian modal Hamiltonian write
$H=S+iB$, where $S$ is real symmetric and $B$ real skew-symmetric.
The real equations are
\begin{equation}
 \dot x=Bx+Sr,\qquad\dot r=Br-Sx.
\end{equation}
Solving the first equation and substituting into the second gives
\begin{equation}
 \dot r(t)=Br(t)-Se^{Bt}x(0)
       -\int_0^tSe^{B(t-s)}Sr(s)\,ds.
\end{equation}
This is an exact reduced equation. It contains the discarded initial
state and a history kernel. Dropping either term defines a different
model unless an approximation and its range of validity are supplied.
In particular, a local damping coefficient cannot be inferred merely
from the fact that a component was projected out.

The same lesson applies to positive-root phase chambers. A projected
real or odd mode can remain on a kinematically selected slice while
its phase exits a positive-root chamber. A phase minimum of an energy
restricted to that slice is not enough to establish stability against
the omitted amplitude and complex-phase directions. The full action
and the resulting projected equation must determine which variations
are physically admitted.

\subsection{An action reduction that keeps the canonical partners}
A controlled finite alternative illustrates how additional inertia
can be derived without postulating a viscous law. Let
\begin{equation}
 L=\frac i2(f^*\dot f-\dot f^*f)-\mathcal E(f),\qquad
 e=(\cos(\theta/2),\sin(\theta/2))^{\mathsf T},
\end{equation}
and write
\begin{equation}
 f=(a+ib)e+cNe,\quad
 N=\begin{pmatrix}0&-1\\1&0\end{pmatrix},\quad b>0.
\end{equation}
The kinetic term is
$b\dot a-a\dot b-bc\dot\theta$. Removing $a$ and $c$ before
variation would remove the canonical partners of amplitude and
phase. Instead set $q=(b,\theta)$, $h=(a,c)$ and
$D_b=\operatorname{diag}(2,b)$. Up to a total derivative,
\begin{equation}
 L=-h^{\mathsf T}D_b\dot q-
       V(q)-j(q)^{\mathsf T}h-\tfrac12h^{\mathsf T}W(q)h.
\end{equation}
For positive definite $W$, the algebraic equation for $h$ is
\begin{equation}
 h=-W^{-1}(D_b\dot q+j),
\end{equation}
and completing the square gives
\begin{equation}
 L_{\rm eff}=\tfrac12(D_b\dot q+j)^{\mathsf T}
                    W^{-1}(D_b\dot q+j)-V.
\end{equation}
The induced kinetic metric is $D_bW^{-1}D_b$ and can depend on
the retained coordinates. Its derivatives belong in the equations
of motion. The hypotheses $W>0$ and $b>0$ are genuine conditions
for this regular reduction.

For the comparison $H=mI+\omega\sigma_y$, with $m>|\omega|$,
the result is
\begin{equation}
 L_{\rm eff}=\dot b^2/m+
       b^2(\dot\theta-2\omega)^2/(4m)-mb^2.
\end{equation}
The angular inertia is $b^2/(2m)$. Its equation
$\ddot\theta=-2(\dot b/b)(\dot\theta-2\omega)$ describes
reversible exchange with the amplitude. Contraction changes the sign
of the term resembling damping. This is why it cannot be interpreted
as positive viscosity without a separately represented dissipative
mechanism. These are modal comparison coordinates, not yet a native
fluid or a physical microscopic clock.

\subsection{Reality of the action and the meaning of a positive mismatch}
Consider ordinary Hermitian internal fields $A_\mu,H_c$, a fixed real
covector $n_\mu$, and
\begin{equation}
 \mathcal B_\mu=iA_\mu+\kappa H_cn_\mu.
\end{equation}
For constant $n$, write
\begin{align}
 D_\mu^AH_c&=\partial_\mu H_c+i[A_\mu,H_c],\\
 F^A_{\mu\nu}&=\partial_\mu A_\nu-\partial_\nu A_\mu
                                      +i[A_\mu,A_\nu],\\
 X_{\mu\nu}&=n_\nu D_\mu^AH_c-n_\mu D_\nu^AH_c.
\end{align}
The connection curvature is $F=iF^A+\kappa X$. Its literal quadratic
trace action and an adjoint-paired trace action differ:
\begin{align}
 \operatorname{Re}\left[-\tfrac14\operatorname{Tr}(FF)\right]
 &=\tfrac14\operatorname{Tr}(F^AF^A)
          -\tfrac{\kappa^2}{4}\operatorname{Tr}(XX),\\
 -\tfrac14\operatorname{Tr}(F^*F)
 &=-\tfrac14\operatorname{Tr}(F^AF^A)
          -\tfrac{\kappa^2}{4}\operatorname{Tr}(XX).
\end{align}
Spacetime indices are contracted in the declared real Lorentz metric.
The traced cross commutator in the second expression vanishes.
These formulas compare different actions, not two notations for the
same variational problem. The choice must be made consistently with
the probe law and its conserved currents.

Even positivity of $H_c$ supplies no general damping theorem. In
the constant-internal, flat-background comparison $H_c=hT_0$, with
$T_0$ Hermitian and $c_T=\operatorname{Tr}T_0^2>0$, the null
curvature term has the form
\begin{equation}
 L_{\rm HC}=\frac{\kappa^2c_T}{2}(n\cdot\partial h)^2,
 \qquad n^2=0.
\end{equation}
Its principal tensor is proportional to $n^\mu n^\nu$ and has
rank one for nonzero null $n$. A positive coefficient in one time
slice is not a four-dimensional Lorentzian wave operator. A physical
propagation or viscosity claim needs the rest of the represented
field system.

\subsection{Induced coframes and a direct same-frame obstruction}
Under a fixed Hermitian spin/internal factorization, let
$j^a=\Psi^*K^a\Psi$ and
$v_H^a=\Psi^*K^aH_c\Psi$. The proposed real coframe can then
have the identity
\begin{equation}
 E^a{}_{\mu}=\tfrac12\partial_\mu j^a+\kappa v_H^a n_\mu.
\end{equation}
If $H_c=0$ and $E$ has rank four, the functions $Y^a=j^a/2$
are local coordinates, and $g=\eta_{ab}dY^a dY^b$ is flat.
The inverse function theorem, rather than a curvature calculation,
establishes this restricted result. A variable adjoint, a geometric
spin connection or a different readout can change the hypotheses.

For the nonzero rank-one update $E=E_0+\kappa v n^{\mathsf T}$,
put $d=1+\kappa n^{\mathsf T}E_0^{-1}v$ and assume $d\ne0$.
The determinant lemma and inverse update give
\begin{equation}
 \det E=(\det E_0)d,\qquad
 n^{\mathsf T}g^{-1}n=
       \frac{n^{\mathsf T}g_0^{-1}n}{d^2},
 \quad g=E^{\mathsf T}\eta E.
\end{equation}
Thus a null covector remains null in this nonsingular update. This
is a geometric compatibility statement. A full solution must also
satisfy the probe, connection, closure and source equations.

There is a stronger consistency test if exactly the same probe,
adjoint and derivative are used both to reconstruct $E$ and to
define the coordinate gamma matrices in the ordered equation. Assume
\begin{equation}
 E^a{}_{\mu}=s_0\operatorname{Re}(\overline\Psi\Gamma^aD_\mu\Psi),
 \quad \Gamma^\mu=(E^{-1})^\mu{}_a\Gamma^a,\quad s_0\ne0,
\end{equation}
with real invertible $E$. Contraction of inverse and coframe yields
\begin{equation}
 4=s_0\operatorname{Re}(\overline\Psi\Gamma^\mu D_\mu\Psi).
\end{equation}
For $R_\Psi=(i\Gamma^\mu D_\mu-\Omega)\Psi$ and
$w=\overline\Psi\Omega\Psi$, it follows that
\begin{equation}
 \operatorname{Im}(\overline\Psi R_\Psi)=4/s_0-\operatorname{Im}w.
\end{equation}
If $w$ is real, the same-frame full-rank reconstruction cannot be
on shell. Reality of $w$ follows, for example, when the mass-like
operator has the required Hermiticity relative to the chosen probe
adjoint. This excludes a particular additional identification; the
Companion has not derived that identification as a native law. It
is not a universal inconsistency theorem for every geometric branch.

\subsection{Finite compatibility versus a field solution}
At a prepared finite jet, normal derivative equations can be written
as a linear system $Nx=r$ after lower data have been fixed. Let
$L_N$ have orthonormal columns spanning $\ker N^{\mathsf T}$.
Then the exact compatibility condition is
\begin{equation}
 L_N^{\mathsf T}r=0.
\end{equation}
For compatible $r$, the general solution is
$x=N^+r+k$, $k\in\ker N$. If a kernel decomposition is supplied,
its coordinates parametrize the remaining normal freedom. For
incompatible $r$, the least-squares value $N^+r$ leaves residual
$-L_NL_N^{\mathsf T}r$ and is not a field solution.

The distinction already appears in
$N=\begin{pmatrix}1&0\\0&1\\0&0\end{pmatrix}$, a three-by-two
matrix with first two rows the identity. A target $(a,b,c)$ is
compatible only when $c=0$, although a least-squares solver returns
$(a,b)$ for every $c$. A small normal residual or a solver success
flag cannot replace the explicit compatibility equation.

The source problem is more informative when higher coordinates split
into adjustable normal variables $w$ and prescribed tangential data $d$:
\begin{equation}
 h_0+Bq+A_fw+A_dd=0.
\end{equation}
Let $L_f$ span the left kernel of $A_f$ and define
\begin{equation}
 J_s=L_f^{\mathsf T}B,\qquad
 b_s(d)=-L_f^{\mathsf T}(h_0+A_dd).
\end{equation}
Eliminating $w$ leaves $J_sq=b_s(d)$. If $J_s$ has full column
rank, the source is unique for compatible data:
\begin{equation}
 (I-J_sJ_s^+)b_s(d)=0,\qquad q(d)=J_s^+b_s(d).
\end{equation}
The remaining higher lift is
$w=-A_f^+(h_0+A_dd+Bq)+k$, $k\in\ker A_f$.
This is a conditional finite source law. Changing which tangential
coordinates are prescribed changes the left-kernel test and can
remove a source constraint previously found by freezing all higher
variables. The derivative mask is therefore part of the mathematical
statement, not a numerical implementation detail.

\subsection{Constraint propagation requires open-set hypotheses}
The declared Cartan comparison uses an invertible real coframe, an
independent Lorentz connection, probe fields and multipliers enforcing
a constitutive coframe relation. Under compatible elimination of the
connection, coordinate and local Lorentz variations give residual
identities. Their source-free specialization requires both the
constitutive residual and the full probe residual to vanish on an
open region. A finite jet satisfying their point values does not
establish those hypotheses.

When the hypotheses do hold, and the spatially projected
trace-reversed coframe evolution equations are imposed, the
gravitational residual admits a symmetric hyperbolic subsidiary
system. With lapse $N_l>0$, shift $\beta$, spatial metric $h$ and
spatial covector $\xi$, its principal symbol can be written
\begin{equation}
 A(\xi)=-(\beta\cdot\xi)I_4+
 N_l\begin{pmatrix}0&(h^{-1}\xi)^{\mathsf T}\\
                         \xi&0_{3\times3}\end{pmatrix}.
\end{equation}
The positive matrix $W=\operatorname{diag}(1,h^{-1})$ symmetrizes
the symbol. Retaining the required lower-order terms and excluding
incoming constraint flux gives an energy estimate
\begin{equation}
 \mathcal E_g'(t)\leq C(t)\mathcal E_g(t).
\end{equation}
Zero initial constraints then remain zero under the stated regularity
conditions. This conditional propagation result does not propagate
the separate constitutive and probe equations from their initial
values. Those field equations are inputs to the subsidiary argument.

The geometric branch therefore has several distinct stages: an
algebraically meaningful coframe, a finite compatible derivative
system, a local field solution, and propagation of its constraints.
The programme has concrete tests for these stages, but a result at
one stage cannot be silently used as a result at the next. The same
discipline that separates positive mass readout from a physical mass
mechanism is needed when a finite source law is proposed as gravity.
\section{\raggedright Discussion and outlook}
The results above locate the main missing information in the transfer
from a common carrier to observables. An HC projector gives an exact
grading. Koide additionally requires a metric-preserving root
representation and a selected norm relation. Positive nonlinear
readouts introduce state-domain restrictions, while geometric
projection introduces an independent autonomy condition. None of
these requirements is fixed by the dimensionality of a representation.

A recent comparison route couples this question to an arithmetically
defined operator. For $\sigma>1$, the logarithmic derivative of zeta
defines the real odd function
\begin{equation}
 a(t)=\operatorname{Im}\frac{\zeta'(\sigma+it)}{\zeta(\sigma+it)}
 =\sum_{n\geq2}\Lambda(n)n^{-\sigma}\sin(t\log n),
\end{equation}
where $\Lambda$ is the von Mangoldt function \cite{DLMF274}.
Multiplication by $a$ anticommutes with parity $R$. Thus the
selfadjoint operator $T_a=-iRA_a$ has symmetric positive and negative
spectral subspaces. A state restricted to either nonzero spectral side
is parity-balanced, including when the spectrum is continuous.
The construction $f(t)=Na(t)^2e^{-t^2}$,
$\psi=(f+i\operatorname{sgn}(a)f)/\sqrt2$ gives a regular normalized
example on the positive side. Its parameter and Fourier variances
are finite. The corresponding Robertson relation is compatible with
the balance, but need not be saturated \cite{Robertson1929}.

This example establishes mathematical compatibility, not physical
selection. The positive spectral restriction, its preparation, the
window and the interpretation of $T_a$ as energy are additional
choices; the parameter is not identified with physical time or position.
The convergent prime series also does not justify substituting
$\sigma=1/2$. Most importantly, the example supplies no map to three
positive lepton roots and does not evade the readout restrictions above.

A physical completion must therefore specify the carrier action,
generation embedding, admissible state domain, observable map and
dynamics on the same represented system. It must then determine which
parameters are predictions and which are measured inputs. Until those
steps are supplied, the framework supports conditional constructions
and consistency tests rather than a recovered particle spectrum or
a unified field theory. The long Companion retains the broader
cosmological and structural proposals with their individual qualifications.
\subsection{An arithmetic comparison operator}

The suggestion that primes, a zeta function, and uncertainty might
participate in the same selection mechanism can be made mathematically
precise without identifying any of them with a physical observable.
This distinction matters: an analogy becomes testable only after its
operators, state space, and domains have been specified. The following
construction is an analytic comparison model. It does not extend the
Companion action or supply a lepton-mass prediction.

Fix $\sigma>1$ and use the positive Hilbert space
$\mathcal H=L^2(\mathbb R,dt)$, with dimensionless $t$. The logarithmic
derivative of the Euler product gives \cite{DLMF274}
\begin{equation}
 -\frac{\zeta'(\sigma+it)}{\zeta(\sigma+it)}
 =\sum_{n\geq2}\Lambda(n)n^{-\sigma}e^{-it\log n},
 \qquad
 \Lambda(p^k)=\log p.
\end{equation}
Here $\Lambda(n)=0$ unless $n$ is a prime power. In particular the
real odd function
\begin{equation}
 a(t)=\operatorname{Im}\frac{\zeta'(\sigma+it)}{\zeta(\sigma+it)}
     =\sum_{n\geq2}\Lambda(n)n^{-\sigma}\sin(t\log n)
\end{equation}
defines a bounded self-adjoint multiplication operator $Af=af$.
Absolute and uniform convergence follow from
$\Lambda(n)\leq\log n$ and convergence of
$\sum (\log n)n^{-\sigma}$. Every fixed derivative is also uniformly
convergent; the relevant majorant has an additional fixed power of
$\log n$. Reflection $Rf(t)=f(-t)$ is a self-adjoint involution and
$RA=-AR$. Thus a parity operator and a genuinely arithmetic odd
partner coexist in one positive Hilbert space.

For a normalized vector, Robertson's inequality \cite{Robertson1929}
becomes
\begin{equation}
 \Delta R\,\Delta A\geq
 \tfrac12|\langle[R,A]\rangle|=|\langle RA\rangle|.
\end{equation}
It does not require $\langle R\rangle=0$. An even normalized vector
has $\langle R\rangle=1$ and both $\langle A\rangle$ and
$\langle RA\rangle$ zero, so it satisfies the inequality without
balance. This is a direct counterexample to an uncertainty-only
selection argument. The inequality restricts a prepared state; it
does not choose the preparation.

\subsection{Why an eigenstate prescription fails}

The bounded operator $T=-iRA$ is self-adjoint and satisfies
$T^2=A^2$ and $RT=-TR$. If a normalized eigenvector with nonzero
eigenvalue existed, anticommutation would force its parity expectation
to vanish. That observation suggests a stronger selection rule than
Robertson alone. In this particular realization, however, there are
no nonzero normalizable eigenvectors, including at eigenvalue zero.

Indeed, the sine series is holomorphic in the strip
$|\operatorname{Im}t|<\sigma-1$ and hence real analytic on the real
axis. Moreover,
\begin{equation}
 a(0)=0,\qquad
 a'(0)=\sum_{n\geq2}\Lambda(n)(\log n)n^{-\sigma}>0.
\end{equation}
Consequently $a^2$ is not constant. An eigenvector of $T$ at real
$\lambda$ would obey $(a(t)^2-\lambda^2)\psi(t)=0$ almost
everywhere. The zero set of that nontrivial real-analytic function
has measure zero, and an $L^2$ vector supported on it is zero.
The argument excludes this eigenstate route, rather than every
possible role for arithmetic. A finite grid could display discrete
eigenvectors while missing this obstruction in the continuum.

\subsection{A regular balanced spectral packet}

One can retain the positive continuous spectral subspace instead.
Let $S=\operatorname{sgn}(A)$; the null set of $a$ is negligible,
so $S^2=I$ and $RS=-SR$. Choose the real even window
$g(t)=e^{-t^2}$, set $f=Na^2g$ with $\|f\|=1$, and define
\begin{equation}
 \psi=\frac{f+iSf}{\sqrt2},\qquad
 \operatorname{sgn}(T)=-iRS.
\end{equation}
Since $f$ is even and $Sf$ odd, their norms agree and their inner
product vanishes. Direct substitution gives
\begin{equation}
 \|\psi\|=1,\quad \langle R\rangle_\psi=0,\quad
 \operatorname{sgn}(T)\psi=\psi,\quad T\psi=|a|\psi.
\end{equation}
The last expression is multiplication by a variable function,
not an eigenvalue equation. Thus it is consistent with the preceding
nonexistence result. If $F_+=\mathbf1_{(0,\infty)}(T)$, then
$F_+f=(f+iSf)/2$ and normalization gives this packet.

The factor $a^2$ also resolves a domain issue. A generic sign change
in $Sf$ could create a jump and an infinite Fourier second moment.
Here $Sf=Na|a|g$, with weak derivatives
\begin{align}
 f'&=N(2aa'g+a^2g'),\\
 (Sf)'&=N(2|a|a'g+a|a|g').
\end{align}
They lie in $L^2$ because $a,a'$ are bounded and the Gaussian and its
derivative are square integrable. There is no delta contribution at a
zero of $a$: the profile vanishes there. In addition $t\psi$ and
$t\psi'$ lie in $L^2$. With $Qf=tf$ and the usual self-adjoint
$P=-i\,d/dt$, both second moments are finite, $QP\psi$ and
$PQ\psi$ exist, and $[Q,P]\psi=i\psi$. Hence
\begin{equation}
 \Delta Q\,\Delta P\geq\tfrac12.
\end{equation}
This is a parameter--Fourier uncertainty relation. Physical position,
momentum, and a dimensional value of $\hbar$ have not been identified.

Balance does not mean uncertainty saturation. Write
$\mu_1=\int |a||f|^2dt$ and $\mu_2=\int a^2|f|^2dt$. For this packet
\begin{equation}
 \langle R\rangle=\langle A\rangle=0,\quad
 \Delta R=1,\quad \Delta A=\sqrt{\mu_2},\quad
 \langle RA\rangle=i\mu_1.
\end{equation}
The Robertson bound is $\sqrt{\mu_2}\geq\mu_1$ and is strict:
equality would require $|a|$ constant almost everywhere on the
support of $f$, whereas its level sets have measure zero. Balance
comes from the chosen positive spectral projection. It is not an
effect of saturating that inequality.

\subsection{Preparation, preservation, and physical identification}

This comparison separates three operations that an effective theory
must not conflate. A spectral projector defines an allowed subspace;
a preparation mechanism produces states in it; a generator preserves
it during evolution. Functional calculus ensures that $e^{-isT}$
commutes with $F_+$, but this does not prepare a generic incoming
state in $\operatorname{im}F_+$, nor identify $T$ with the native
Hamiltonian. The window, its scale, and the suppression of zeros
are additional choices. None follows from the Euler product.

The construction is confined to $\sigma>1$. Continuing the scalar
zeta function does not automatically continue this bounded operator
or its domain. On the critical line, poles of the logarithmic
derivative require a separate analysis. No statement here proves
the Riemann hypothesis or turns its zeros into the observed masses.
The useful outcome is narrower: an arithmetic odd partner and a
regular balanced packet are compatible with uncertainty, while
their physical preparation and mass readout remain open.

\subsection{What would constitute progress}

For the mass branch, a decisive next result would specify the native
operator $\widetilde R$, its compressed transition $B$, the positive
three-dimensional readout, and the rule selecting its phase. A fit
to three masses tests the flexibility of a chosen parameterization;
it does not supply these missing definitions. A scale-independent
ratio can constrain a mechanism before an overall mass scale is
known, but agreement of a ratio cannot identify that scale or the
underlying dynamics.

For the generation branch, an admissible seed map must retain rank
three and satisfy the grading requirements through the actual
operator chain. Mixing requires a second physically identified
operator in the same comparison space. For the geometric branch,
an invertible coframe, compatible finite derivative data, and a
local solution are separate obligations. Constraint propagation is
valuable after its field-equation hypotheses are met.

The database edition makes these distinctions inspectable. It
stores statement versions, dependencies, source hashes, and the
editorial selection used to assemble this paper. Its purpose is to
prevent an algebraic possibility, a conditional calculation, and
a physical identification from becoming indistinguishable during
rewriting. The present work supplies explicit constructions and
obstructions for these interfaces. A completed phenomenological
theory would additionally need independently specified observables,
units, preparation rules, and comparison with measured data.
\begin{thebibliography}{99}
\bibitem{Baez2002}
J.~C.~Baez, \emph{The Octonions}, Bulletin of the American Mathematical
Society \textbf{39} (2002), 145--205,
doi:10.1090/S0273-0979-01-00934-X.
\bibitem{Robertson1929}
H.~P.~Robertson, \emph{The Uncertainty Principle}, Physical Review
\textbf{34} (1929), 163--164, doi:10.1103/PhysRev.34.163.
\bibitem{QSC411}
M.~Brendecke, \emph{Quantum Sphaera Companion: A Structure-First Mathematical
Unfolding}, v4.11.0 (2026), doi:10.5281/zenodo.24108818.
\bibitem{DLMF274}
NIST Digital Library of Mathematical Functions, Section 27.4,
\emph{Euler Products and Dirichlet Series}, Eq.~27.4.12,
dlmf.nist.gov/27.4 (accessed 11 November 2026).
\end{thebibliography}
\clearpage
\appendix
\section{Supporting mass-readout arguments}
For completeness, the original mathematical arguments for the conditional mass-readout statements are reproduced below. Their hypotheses refer to the comparison domains stated in the text; they do not identify a native physical readout.
Let $x_i=\sqrt{m_i}\geq0$, $x\ne0$, and write
\[
 A=\frac{x_0+x_1+x_2}{3}>0,\qquad
 u=A(1,1,1),\quad d=x-u,\quad \sum_i d_i=0.
\]
The cyclic permutation preserves the one-dimensional symmetric line
and its two-dimensional orthogonal plane. It preserves their norms
separately; it does not prescribe their ratio.
\begin{proposition}[Koide as an equal norm condition]
\label{v480:koide-equal-norm}
For these nonnegative principal roots,
\[
 K=\frac{\sum_i x_i^2}{(\sum_i x_i)^2}
   =\frac13+\frac{\|d\|^2}{9A^2}.
\]
Consequently $K=2/3$ if and only if
$\|d\|^2=\|u\|^2=3A^2$.
\end{proposition}
\begin{proof}
Since $u\perp d$, $\|x\|^2=3A^2+\|d\|^2$, while
$(\sum_i x_i)^2=9A^2$. Division proves both assertions.
\end{proof}
\begin{proposition}[A sufficient HC shaped mass root realization]
\label{v480:koide-hc-mixing}
In a declared real three-root representation, let
$e=(1,1,1)/\sqrt3$, let $J_m$ be an orthogonal involution, and put
$C_m=(I+J_m)/2$. Suppose $x=\alpha C_m e$ is nonzero and
componentwise nonnegative, with $\alpha>0$. For
$h=\langle e,J_m e\rangle$, one has $h>-1$ and
\[
 K=\frac{2}{3(1+h)}.
\]
Hence Koide holds precisely when $\langle e,J_m e\rangle=0$.
\end{proposition}
\begin{proof}
An orthogonal involution is self-adjoint, so $C_m$ is an
orthogonal projector. Thus
$\|C_m e\|^2=\langle e,C_m e\rangle=(1+h)/2>0$ and
$(\sum_i x_i)^2=3\alpha^2(1+h)^2/4$. Dividing
$\|x\|^2=\alpha^2(1+h)/2$ by this denominator gives the formula.
\end{proof}
\begin{proposition}[Radius and phase in the three root parametrization]
\label{v480:koide-radius-phase}
Every nonzero nonnegative three-root vector admits
\[
 x_k=A\bigl(1+r\cos(\theta+2\pi k/3)\bigr),\qquad r\geq0,
\]
with $K=1/3+r^2/6$. The condition $K=2/3$ is precisely $r=\sqrt2$;
it does not fix $A$ or $\theta$.
\end{proposition}
\begin{proof}
The vectors $(\cos(2\pi k/3))_{k=0}^2$ and
$(\sin(2\pi k/3))_{k=0}^2$ form an orthogonal basis of the
zero-sum plane, each with squared norm $3/2$. Expressing $d/A$
in that basis gives coefficients $r\cos\theta$ and
$-r\sin\theta$. Therefore $\|d\|^2=(3/2)A^2r^2$ and the
previous proposition gives the formula for $K$. If $r=0$ the
phase is arbitrary.
\end{proof}
\begin{proposition}[Cyclic equivariant linear maps preserve symmetric inputs]
\label{v480:koide-equivariant-obstruction}
Let $P$ cyclically permute three real coordinates. If a linear map
$T$ satisfies $TP=PT$, then $T(A(1,1,1))$ remains on the symmetric
line. It cannot produce a nondegenerate root hierarchy from that input.
\end{proposition}
\begin{proof}
$Pu=u$ implies $P(Tu)=T(Pu)=Tu$. The fixed space of $P$ is
exactly the line spanned by $(1,1,1)$.
\end{proof}
\begin{proposition}[descent]
\label{v480:hc-mass-descent}
Let $H$ be a real Hilbert space, $J$ an orthogonal involution, and
$B:H\to\mathbb R^3$ a bounded coisometry, $BB^*=I_3$. If
$\ker B$ is $J$-invariant, then
\[
 J_m=BJB^*,\qquad J_m^*=J_m,\qquad J_m^2=I_3,
 \qquad BC=C_mB,\quad C_m=(I_3+J_m)/2.
\]
Conversely, any relation $BJ=J_mB$ requires $\ker B$ to be
$J$-invariant. The descended operator is unique because $B$ is onto.
\end{proposition}
\begin{proof}
Put $E=B^*B$, the orthogonal projection onto $(\ker B)^\perp$.
Since $J$ is self-adjoint and preserves $\ker B$, it also preserves
its orthogonal complement, so $JE=EJ$. Thus $J_m^*=J_m$ and
$J_m^2=BJEJB^*=BEB^*=I_3$. Also
$J_mB=BJE=BJ$, since $J$ preserves $\ker B$ and $B$ kills it.
The identity for $C$ follows by averaging. Conversely,
$Bk=0$ implies $BJk=J_mBk=0$. Surjectivity gives uniqueness.
\end{proof}
\begin{proposition}[odd obstruction]
\label{v480:hc-mass-odd-obstruction}
If $Jf=-f$, then $Cf=0$. Consequently any homogeneous linear
root readout $x=BCf$ is zero. Any readout $F(Cf)$ with $F(0)=0$
has the same obstruction on this input lane. If instead a seed is
$J$-fixed and the descended linear route has $Bf_0=e$, then
$J_me=e$ and its nonzero root output has $K=1/3$.
\end{proposition}
\begin{proof}
The first statement follows from $Cf=(f+Jf)/2=0$.
For the second, $J_me=BJf_0=Bf_0=e$, so $C_me=e$ and
all three roots are equal.
\end{proof}
\begin{proposition}[equivariant obstruction]
\label{v480:hc-mass-equivariant-obstruction}
Let $U$ preserve an admitted state set and let $P$ cyclically permute
the three root coordinates. For any readout $F$ satisfying
$F(Uf)=PF(f)$, an input $Uf_0=f_0$ has
$F(f_0)=a(1,1,1)$. If this is a nonzero nonnegative root vector,
then $K=1/3$. No linearity assumption on $F$ is needed.
\end{proposition}
\begin{proof}
Equivariance gives $PF(f_0)=F(Uf_0)=F(f_0)$. The fixed vectors
of the three-cycle are exactly the constant triples.
\end{proof}
\begin{proposition}[triole bound]
\label{v480:hc-mass-triole-bound}
Let $z_0+z_1+z_2=0$ in a normed vector space, not all $z_i=0$.
Under the specific identification $m_i=\kappa\|z_i\|^2$ with a
common $\kappa>0$, the principal-root ratio satisfies
\[
 \frac13\leq K\leq\frac12.
\]
In particular it cannot equal $2/3$.
\end{proposition}
\begin{proof}
Put $r_i=\|z_i\|$ and $s=\sum_i r_i>0$. The triangle inequality
gives $r_i\leq s-r_i$, hence $r_i\leq s/2$. Therefore
$\sum_i r_i^2\leq(s/2)\sum_i r_i=s^2/2$.
Cauchy--Schwarz gives the lower bound, and the common mass scale
cancels from $K$.
\end{proof}
\begin{proposition}[fit freedom]
\label{v480:hc-mass-fit-freedom}
For every nonzero nonnegative triple $x$, there exist an orthogonal
involution $J_x$ and $\alpha_x>0$ such that
$x=\alpha_x(I+J_x)e/2$. Thus existence of an HC-shaped operator
chosen after the masses is not a prediction of Koide or a hierarchy.
\end{proposition}
\begin{proof}
Set $C_x=xx^{\mathsf T}/\|x\|^2$ and $J_x=2C_x-I_3$.
This is an orthogonal involution. Since $s=\sum_i x_i>0$,
$C_xe=sx/(\sqrt3\|x\|^2)$, so choose
$\alpha_x=\sqrt3\|x\|^2/s$.
\end{proof}
\end{document}
